\section*{Abstract}\label{section:Abstract}

Diese Arbeit entwickelt und evaluiert ein webbasiertes Validierungssystem, das Bürger bei der räumlich-semantisch korrekten Platzierung von Objekten in einer Augmented-Reality-Partizipationsplattform für die Stadtplanung unterstützt. Auf Basis eines systematischen Requirements-Engineerings mit einem Experteninterview wird zunächst die \textit{ARVal Constraint Syntax} (ACS) entworfen und formal als JSON-Schema beschrieben. Auf Grundlage des Schemas wird ein mehrschichtiges Softwaresystem für die Regelprüfung in WebXR implementiert. Zur Sicherstellung der Laufzeitfähigkeit werden Performance-Profiling und Optimierungen (u. a. Geometry-Batching, Memoisation) durchgeführt.

Die Nutzbarkeit wird in einem kontrollierten Vergleichstest ohne und mit aktiviertem System untersucht. Mit Validierung sinkt der mittlere NASA-TLX-Wert zur mentalen Belastung signifikant von 39 auf 18 Punkte und die Aufgabenzeit von 178s auf 111s. Der SUS-Wert zur Nutzbarkeit des Systems steigt von 65 auf 80 Punkte. Qualitative Think-Aloud-Daten bestätigen, dass verständliche Fehlermeldungen und visuelle Markierungen das Regelkonzept schnell internalisieren lassen.

Die Arbeit zeigt damit, dass ein regelbasiertes Feedback in Augmented Reality (AR) nicht nur technische Machbarkeit, sondern auch einen klaren Mehrwert für Intuitivität und Effizienz bietet. ACS, das offene JavaScript-Framework und die empirischen Ergebnisse bilden eine Grundlage für künftige Forschung zu komplexeren Regelwerken und großskaligen Feldstudien in der partizipativen Stadtplanung.